\scanpage{101}
\mspage{45}
% archive sheet number: 100
\unsure{Pour|Par} la suite, \unsure{par} un ``localiseur fondamental'', nous entendrons un ens.

\[
 W\subset\mathrm{Fl}\,\mathrm{Cat}
\]

satisfaisant \inserted{(tout au moins)} les axiomes

\[
 \mathrm L1,\ \mathrm L3\ \text{quater},\ \mathrm L4\ \text{bis},\ \mathrm L7
\]

i.e. \deleted{\unsure{satisferaient les}}

\noindent\begin{minipage}{.035\linewidth}
\begin{tikzpicture}
\draw (0.25,1.45) .. controls (0.12,1.28) and (.29,.14) .. (0,0) .. controls (.29,-.08) and (.12,-1.28) .. (.25,-1.45);
\end{tikzpicture}
\end{minipage}\begin{minipage}{.93\linewidth}
a) \uline{Saturation} de $W$ (cond. a) b) c) p.~1)

b) $W_\infty\subset W[\subset W_0]$\inserted{cf d)}

c) Si $f:X\longrightarrow Y$ sur $S$ est loc. sur $S$ une $W$-équiv., alors $f\in W$ \inserted{(L4 bis)}

d) Pour que $f:X\longrightarrow Y$ soit dans $W$, il f. et s. que $\pi_0(f)$ soit bijectif, et que les $f_i:X_i\longrightarrow Y_i$ induits sur les compos. conn. soient dans $W$.
\end{minipage}

% slanted
\margin{Je ne suis plus sûr s'il est vraiment utile d'imposer $W_\infty\subset W$. \deleted{\illegible{1}} Prendre au \deleted{lieu de} ça l'axiome de l'objet final, \unsure{soit b')|L iv bis b')}. Si on veut a) b') c) + L5, alors on aura $W\subset\unsure{W_\infty}$.}

% slanted
\margin{\uline{NB} a) b) c) + L5 implique $W=W_\infty$, cf annotation page 43.}

\deleted{Par contre, les \unsure{axi}} Ces axiomes sont vérifiés par tous les $W$ que je connais, et qu'on aurait envie de regarder comme localiseurs. Par contre, L5, L6, L6 bis (p.~43,44) sont des axiomes de nature beaucoup plus spéciale, j'ai envie de dire qu'ils distinguent les localiseurs de nature proprement ``\uline{homotopique}'' de ceux qui sont de nature ``\uline{homologique}'' (ou ``cohomologique''). Il faudrait que je m'assure qu'ils sont satisfaits pour les $W_i$ ($i$-équiv. d'homotopie).

% slanted
\margin{C'est \uline{faux}, cf annotation ci-dessus}

Le fait que L4 bis \inserted{i.e c) ci-dessus} implique l'axiome dual résulte du fait \deleted{que si $f:X\longrightarrow Y$, $f\in W$} que L4 bis implique ceci

\uline{Proposition 2} Soit $f:X\longrightarrow Y$ dans $\mathrm{Cat}$, alors $f\in W$ ssi $f^o:X^o\mathbin{\unsure{-|\to}}Y^o$ est dans $W$.

Se voit par l'argument de I, page 75 Corr.~2) -- n'utilise que L4, pas $=$ L1, L3.

